\section{Random formulas in First-Order Logic}


\subsection{Random formulas in First-Order Logic}

Let there be a random distribution of models in First-Order Logic with a fixed domain of size $\inddim$.
A random formula of order $\indorder$ is then the random substitution tensor
\begin{align*}
    \folexformulaat{\indvariableof{[\indorder]}} \, .
\end{align*}

Properties:
\begin{itemize}
    \item Always: The marginal distribution of each coordinate is a Bernoulli distribution with parameter $p$.
    \item Restrictive: We sometimes assume the coordinates to be i.i.d.
\end{itemize}


\subsubsection{Conjunctions of random formulas}

The conjunction of random formulas is a random tensor network of them.


\begin{example}[SPARQL queries as random formulas]
    Consider a SPARQL query with a fixed number of triple patterns.
    Each triple pattern is a random formula of order $3$.
    The conjunction of the triple patterns is then a random tensor network of the random formulas.
\end{example}

When each predicate is an independent random formula of i.i.d. Bernoulli coordinates, we have a $2$-Orlicz random tensor network (see \cite{goessmann_uniform_2021}).




\subsection{Statistical models of knowledge graphs}

Setup:
\begin{itemize}
    \item Atom extraction queries: Treated as deterministic
    \item Importance query: I.i.d. bernoulli coordinates (i.e. $2$-Orlicz). Rescaled for isotropy.
\end{itemize}

Aim:
\begin{itemize}
    \item Predict the average of specific functions. This is the expected contraction with the isotropic importance query.
\end{itemize}

\subsection{Concentration of single expectations}

We first show concentration of the expectation to single formulas
Bernoulli coordinates are bounded and therefore sub-Gaussian, that is $2$-Orlicz.
We can therefore apply the results on Orlicz tensor networks to show moment bounds.

\subsection{Uniform concentration of multiple expectations}

Having moment bounds for single expectations, we use chaining with Mendelsons $\Gamma$ functionals to show uniform concentrations.
We can then show probabilistic guarantee recovery on least squares regression.

